\section{Local remainders and the interacting vacuum}
\label{rel4-section}

This section separates an elementary useful remainder estimate from
the additional vacuum coercivity required for a dynamical low-energy
theorem. The on-site sea is not the interacting vacuum, even after a
nontrivial electric-flux compression. All statements concern the
unimproved Wilson Hamiltonian \eqref{lat-hamiltonian}, with $n\ge3$,
$N=n^3$, $m>0$, and $g^2>0$.

\subsection{The bare excitation number cannot have a vacuum-relative gap}

Choose the $\beta$ eigenbasis and fill the two negative-$\beta$ modes
at every site. Denote that normalized fermionic vector by $\Omega_b$.
Let $N_e$ count occupied positive-$\beta$ modes and let $N_p$ count
holes in the negative-$\beta$ modes. Then
\begin{equation}
 N_{\rm exc}=N_e+N_p=d\Gamma(\beta)+2N,\qquad
 Q=N_e-N_p,\qquad
 \ker N_{\rm exc}=L^2(U(1)^{\mathcal E})\otimes\C\Omega_b .
 \label{rel4-excitation-definition}
\end{equation}
These operators commute with every local Gauss transformation.
They do not commute with the full Hamiltonian.

\begin{proposition}[Bare-sea obstruction, also at bounded flux]
\label{rel4-bare-sea}
No eigenvector of $H$ belongs to $\ker N_{\rm exc}$. Consequently any
normalized physical vacuum ground state $\Psi_0$ satisfies
$\langle\Psi_0,N_{\rm exc}\Psi_0\rangle>0$.
The same assertion holds for the finite-dimensional compression
$P_FHP_F$ to $|E_e|\le F$ on every link, for every integer $F\ge1$.
In particular, a bound
\begin{equation}
 \langle N_{\rm exc}\rangle_\Psi
 \le \frac{C}{m}\bigl(\langle H\rangle_\Psi-E_0+C_s(Z_{\rm tot}+K)\bigr)
 \label{rel4-false-coercivity}
\end{equation}
cannot hold on all sufficiently low-energy states, even at zero
source, for finite constants $C$ and $C_s$; $C_s$ has energy units.
\end{proposition}
\begin{proof}
Write $V=d\Gamma(\sum_j\alpha_jv_j)$. All other terms of $H$
preserve $\ker N_{\rm exc}$. The Wilson hopping within the filled
negative block vanishes on $\Omega_b$, and the positive block is empty.
For each positively oriented link $e=(x,j)$, $V$ creates an electron
at one endpoint and a hole at the other. The two orientations give
orthogonal fermionic states. For $n\ge3$, different oriented links
also give orthogonal states. The spin matrix is $\sigma_j/(2ia)$
up to a sign and a unit-modulus link factor, with
$\operatorname{tr}(\sigma_j^*\sigma_j)=2$. Hence for any link
wavefunction $f$,
\begin{equation}
 \|V(f\Omega_b)\|^2=\frac{3N}{a^2}\|f\|^2 .
 \label{rel4-sea-pair-norm}
\end{equation}
This vector is orthogonal to $\ker N_{\rm exc}$. The eigenvalue
equation therefore cannot hold on a nonzero $f\Omega_b$.
The compact-resolvent realization guarantees a ground state in the
physical vacuum sector.

For the compressed statement put $P_F=\prod_e\mathbf1_{[-F,F]}(E_e)$.
The projection commutes with $N_{\rm exc}$ and with Gauss's law.
On the one-link range of $P_F$, the two compressed shifts obey
\[
 \|P_FUf\|^2+\|P_FU^*f\|^2
 =2\|f\|^2-\|\mathbf1_{\{E=F\}}f\|^2
              -\|\mathbf1_{\{E=-F\}}f\|^2
 \ge\|f\|^2 .
\]
The last inequality uses $F\ge1$; the two endpoint projections are
then orthogonal. With all other link projections present, the same
identity holds by tensoring with their ranges. The preceding
orthogonality computation now gives
\[
 \|P_FVP_F(f\Omega_b)\|^2\ge\frac{3N}{2a^2}\|f\|^2 .
\]
Thus the compressed eigenvalue equation also cannot hold on the
bare-sea subspace. At $K=Z_{\rm tot}=0$, a ground state of the
original Hamiltonian makes the right side of
\eqref{rel4-false-coercivity} zero and the left side strictly positive.
For the compressed assertion, replace $H$ and $E_0$ there by
$H_F=P_FHP_F$ and its physical vacuum energy $E_{0,F}$.
Its own ground state gives the same contradiction.
\end{proof}

This is a counterexample to the specified \emph{bare} number
coercivity, not to a gap for an appropriately dressed heavy-particle
projection. The case $F=0$ deletes all derivative hopping and is a
different, degenerate compression.

\subsection{A uniform bound on the virtual-pair density}

There is a useful quantitative replacement, which displays rather
than removes the vacuum contribution. Put
\[
 \mu=m+\frac{3r}{a},\qquad H_g=H_E+H_B,\qquad
 E_g=\inf\sigma(H_g),\qquad E_{\rm ref}=E_g-2\mu N .
\]
The pure gauge operator has a unique strictly positive ground
function: its positive torus heat kernel and bounded real potential
give a positivity-improving semigroup. Gauge invariance and uniqueness
make that function gauge invariant. Its product with $\Omega_b$ is
a physical vacuum trial state, so $E_0\le E_{\rm ref}$.

\begin{proposition}[Excitation density with an explicit volume term]
\label{rel4-density}
For every $0<\eta<m$, as quadratic forms,
\begin{equation}
 H\ge E_{\rm ref}+(m-\eta)N_{\rm exc}
                  -\frac{3N}{2\eta a^2}.
 \label{rel4-density-form}
\end{equation}
In particular, any normalized state with
$\langle H\rangle\le E_0+\Delta$, $\Delta\ge0$, satisfies
\begin{equation}
 \langle N_{\rm exc}\rangle
 \le \frac{2\Delta}{m}+\frac{6N}{(ma)^2}.
 \label{rel4-density-energy}
\end{equation}
No source restriction or weak gauge coupling is needed for this
estimate; the statement can be restricted to any physical sector.
\end{proposition}
\begin{proof}
The pair-preserving part $H_d=H_g+d\Gamma(\beta M(U))$ obeys
\[
 H_d=H_g-2\operatorname{Tr}_{\ell^2(X)}M(U)+T_e(U)+T_p(U)
 \ge E_{\rm ref}+mN_{\rm exc}.
\]
Here $T_e\ge mN_e$, $T_p\ge mN_p$, and
$\operatorname{Tr}M=\mu N$ exactly. The transpose in the hole
operator preserves positivity.

At a fixed link configuration the pair-changing block is
$S(U)=\sum_j\sigma_jv_j(U)$. Its singular-value decomposition is
used only for a pointwise fermionic inequality, not to conjugate the
electric kinetic energy. For one singular value $s$, the operator
$s(e^\dagger p^\dagger+pe)$ acts on the even two-mode subspace as
$\left(\begin{smallmatrix}0&s\\s&0\end{smallmatrix}\right)$.
Consequently
\[
 \pm s(e^\dagger p^\dagger+pe)
 \le\eta(n_e+n_p)+\sqrt{\eta^2+s^2}-\eta
 \le\eta(n_e+n_p)+\frac{s^2}{2\eta}.
\]
The odd subspace has zero pair matrix and satisfies the same inequality.
Summing singular modes gives
\[
 \pm V(U)\le\eta N_{\rm exc}
             +\frac{\operatorname{Tr}S(U)^2}{2\eta}.
\]
The Pauli trace removes cross terms. For $n\ge3$,
$\operatorname{Tr}T_j^2=0=\operatorname{Tr}(T_j^*)^2$, and hence
\[
 \operatorname{Tr}S^2
 =2\sum_j\operatorname{Tr}v_j^2=\frac{3N}{a^2}.
\]
The resulting pointwise bound integrates to a form bound on the
full Hilbert space. Adding the estimate on $H_d$ proves
\eqref{rel4-density-form}. Use $\eta=m/2$,
$E_0\le E_{\rm ref}$, and the assumed energy budget to obtain
\eqref{rel4-density-energy}.
\end{proof}

Thus the heavy regime controls a density of bare virtual pairs.
It does not turn their total number into $O(K)$ at fixed $ma$ and
unbounded volume.

\subsection{What annihilation of a local sea actually proves}

\begin{lemma}[Local-annihilation remainder estimate]
\label{rel4-annihilation}
Let $P_x$ project onto the on-site bare sea and put
$P_X=\prod_{x\in X}P_x$, tensoring with the identity on gauge variables.
Let bounded self-adjoint local terms $R_X$ satisfy $R_XP_X=0$.
If
\[
 \varepsilon=\sup_x\sum_{X\ni x}\|R_X\|<\infty,
 \qquad R=\sum_XR_X ,
\]
then $-\varepsilon N_{\rm exc}\le R\le\varepsilon N_{\rm exc}$.
\end{lemma}
\begin{proof}
Self-adjointness gives $P_XR_X=0$, so
$R_X=(I-P_X)R_X(I-P_X)$ and
$-\|R_X\|(I-P_X)\le R_X\le\|R_X\|(I-P_X)$.
The commuting on-site projections satisfy
$I-P_X\le\sum_{x\in X}(I-P_x)\le\sum_{x\in X}n_{{\rm exc},x}$.
Summing proves the assertion.
\end{proof}

The lemma validates the algebraic part of the proposed design
principle. Combined with \eqref{rel4-density-energy}, however, it
still leaves the error
$\varepsilon[2\Delta/m+6N/(ma)^2]$.
If instead $R_XP_X=A_X(U)P_X$ for a nonconstant gauge operator,
subtracting a number does not establish the lemma's hypothesis.
The operator $A_X(U)$ must be retained in the effective gauge
Hamiltonian or controlled in differences of sourced and unsourced
expectations. Its being the same operator in both problems does
not make its two ground-state expectations equal.

A genuinely dressed version would replace $P_x$ and $N_{\rm exc}$
by precisely constructed projections and an excitation number for
which the interacting vacuum has zero expectation, and would prove
an energy bound for that number. Finite-order normal ordering does
not by itself establish either property. In particular, an exact
annihilation claim at $K=0$ must be checked on the actual interacting
vacuum, not just on the reference product vector.

\subsection{The missing dynamical Schur-complement estimate}

For clarity, the formal reduction can be stated as an exact
conditional lemma. Let $P$ be a specified retained projection in a
physical sector, $Q=I-P$, and write
$H=\left(\begin{smallmatrix}A&B\\B^*&D\end{smallmatrix}\right)$.
If $D-z\ge\delta Q>0$ as forms, then completion of the square gives
\begin{align}
 \langle(u,v),(H-z)(u,v)\rangle
 &=\langle u,[A-z-B(D-z)^{-1}B^*]u\rangle \nonumber\\
 &\quad+\|(D-z)^{1/2}[v+(D-z)^{-1}B^*u]\|^2 .
 \label{rel4-schur-identity}
\end{align}
This holds on the form domain when the displayed compositions are
defined; in particular it is an ordinary matrix identity after a
finite electric compression. If, in addition,
$B(D-z)^{-1}B^*\le\varepsilon N_{\rm eff}$ on $P$, one obtains
the lower Schur bound $A-z-\varepsilon N_{\rm eff}$.
On a no-positron retained space with $N_{\rm eff}=Z_{\rm tot}$,
this is genuinely proportional to source charge.

The required lower bound on the \emph{full} pair block $D-z$ and
the source-relative bound on its Schur correction have not been
proved here. The elementary flux-restricted denominators in
Section~\ref{flux4-section} do not imply these operator hypotheses.
One must also control the retained gauge Hamiltonian and its
vacuum normalization. A charged-sector gap concerns different
values of $Q$ and supplies no such estimate for neutral pairs.

The local Schrieffer--Wolff theorem of
Bravyi--DiVincenzo--Loss~\cite{BDL}, Theorem~3, has a remainder
$O(N|\epsilon|^{k+1})$ for its specified finite-dimensional,
on-site-gapped setting. Its locality theorem does not replace the
two additional estimates above. The relative-form and isolated-band
assumptions in extensions to unbounded systems must likewise be
checked rather than imported by name.
